\scanpage{74}

% top, line 1
L'iso can $\mathrm{hex}\simeq\mathrm{tria}$ peut s'expliciter ainsi :

% top, line 2
Soit $h=$

% top, line 2, centre
\diagram{\begin{tikzcd}[ampersand replacement=\&,row sep=small,column sep=small]
\& a\arrow[r,dash]\arrow[dl,dash]\& b'\arrow[dr,dash]\&\\
c'\arrow[dr,dash]\&\&\&c\arrow[dl,dash]\\
\&b\arrow[r,dash]\&a'\&
\end{tikzcd}}

% upper, lines 2–3
un hexagone, \deleted{alors} formé de deux trièdres parallèles $t=\{a,b,c\}$, $t'\unsure{=}\{a',b',c'\}$.

% left margin, bottom to top
% margin note written sideways
\margin{Autre procédé : pour deux côtés opposés, le troisième côté du triangle qu'elles déterminent est le $=$. On a ainsi $3$ pts (correspondants \deleted{aux} aux trois diagonales $=$ \unsure{codiagonales|diagonales} de l'hexagone) formant un trièdre \ldots}

% upper-middle, lines 4–12
Alors $\exists!$ pavé $P\supset h$ (c'est aussi l'unique pavé contenant $\{a,b,c\}$, \unsure{--|ou} $\{a',b',c'\}$), et $t''=P-h=\{a'',b'',c''\}$ est un trièdre par. aux deux précédents (et c'est l'unique tels trièdres, distincts des précédents). Ses él. sont en corr. $1-1$ avec les diagonales de $h$, à la diagonale $\{a,a'\}$ correspondant l'\inserted{unique} él. $a''$ \deleted{tel que} \inserted{(tq)} \unsure{non|bien} lié à $a,a'$, donc lié à $b,b',c,c'$. C'est aussi l'unique el. de $S$ ayant ces propriétés,

% middle to lower, lines 13–17
(car c'est le troisième sommet du triangle de côté $\{b,c'\}$, \unsure{--|-- ou} du triangle de côté $\{b',c\}$ -- ces deux sommets se trouvent \inserted{devoir} être égaux (facile à voir directement\ldots\ldots\ \deleted{\illegible{1}}) et \deleted{il est évident que} \inserted{être liés à $a,a'$} \deleted{directement} (évident directement également.)

% lower, lines 18–20
Inversement, si $t''=\{a'',b'',c''\}$ est un trièdre, elle est contenue dans un unique pavé $P$, et $P-t''$ est un hexagone.

% foot, lines 21–23
Question : Soit $h$ un \inserted{(petit)} hexagone, correspondant à un trièdre $t$ (qui le détermine). Notons que $t$ détermine \uline{également} un (grand) hexagone.

% archive sheet number: 73
